\chapter{Files and directories}
\label{ch:files}

Phase~5 ended with a mount that reached \code{core/} and refused: the on-disk
format had a root tree and an extent tree and no third thing, so there was no
root directory to build an inode from. This chapter is the third tree, and what
it takes to read a file out of it.

The work divides oddly. Three \emph{format} decisions --- how a directory entry
is keyed, what an inode records, how a file's bytes are described --- carry
almost all the weight, because each is unchangeable once \fnname{mkfs} has
written one. The kernel code that consumes them is mostly borrowed: the page
cache, the dentry cache and the \code{getdents} machinery are the VFS's, and
what bitterfs supplies is a few hundred lines that answer questions about keys.

\section{The FS tree}

One tree holds the entire namespace. Five item types live in it, and the
key-type numbering in \code{format/bitterfs\_format.h} \emph{is} their sort
order:

\begin{quote}
\itemtype{INODE\_ITEM} 1, \itemtype{INODE\_REF} 25, \itemtype{DIR\_ITEM} 48,
\itemtype{DIR\_INDEX} 72, \itemtype{EXTENT\_DATA} 96
\end{quote}

Keys sort \keytriple{objectid}{type}{offset}, so everything about one object is
contiguous, its inode first and its data map last. Two properties follow, and
both were designed for rather than discovered.

\paragraph{An object's items are a range.} Deleting a file is a scan from
\keytriple{N}{0}{0} to \keytriple{N}{255}{-1}, which is what
\code{BITTER\_KEY\_TYPE\_MAX} exists to bound. Truncating one is the same scan
restricted to \itemtype{EXTENT\_DATA}, and it is contiguous because that type
sorts last.

\paragraph{A file and a directory differ only in which types appear.} Both are
an objectid with an \itemtype{INODE\_ITEM}. A file has \itemtype{EXTENT\_DATA}
items; a directory has \itemtype{DIR\_ITEM} and \itemtype{DIR\_INDEX} ones. The
tree does not know or care which --- \field{mode} is what says. An FS tree
containing nothing but \keytriple{256}{INODE\_ITEM}{0} is already a complete,
valid, empty root directory, which is exactly what \fnname{mkfs} wrote first.

\section{Two items per directory entry}

Every entry is stored \emph{twice}, with identical payloads under two keys:

\begin{quote}
\keytriple{dir}{DIR\_ITEM}{hash(name)} \quad for lookup\\
\keytriple{dir}{DIR\_INDEX}{seq} \quad for \fnname{readdir}
\end{quote}

Because two operations need opposite orderings and no single one serves both.
Lookup has a name and wants one search, so the key must sort by name.
\fnname{readdir} must be \emph{resumable} --- \code{getdents} hands userspace a
cookie and continues later --- so it needs a stable monotonic ordering, which a
hash cannot give: hashes collide, so the cookie would not be unique, and a
newly created entry can land \emph{behind} the current position and never be
returned.

They are not interleaved. Type dominates offset, so a directory is one run of
hash entries followed by one run of index entries, in different orders, and each
operation walks only its own. A directory with $n$ entries is $1 + 2n$ items;
create writes two, unlink removes two, rename touches four, and nothing yet
checks that the two runs agree.

\subsection{The index is a cookie, not a count}

\field{next\_dir\_index} lives in the directory's own inode item and only ever
rises. A directory that has had a thousand files created and deleted holds 1002
and no entries.

Monotonicity is the whole correctness argument: \fnname{readdir}'s
\code{f\_pos} \emph{is} this number, so reusing a freed value would let a
resumed enumeration miss an entry it never saw, or return one it already had.
Gaps cost nothing --- a search for a deleted index returns the next real entry
--- which is the same property that makes deletion-riddled directories walkable.

It starts at 2, because \fnname{dir\_emit\_dots} consumes positions 0 and 1 for
\code{.} and \code{..}. Those two are synthesised from the dentry tree and
stored nowhere, which is also why the root directory has \field{nlink} 2 with an
empty tree.

\subsection{The hash, and the collision it guarantees}

\itemtype{DIR\_ITEM}'s offset is CRC-32/ISCSI over the name bytes --- the same
algorithm and the same parameters \code{bt\_block\_csum} uses for blocks, so a
key can be checked by hand against any off-the-shelf implementation when a
lookup misbehaves and the question is whether the key or the search is wrong.
It is specified in the format header rather than left implicit in
\code{core/dir.c}, because it is as much part of the on-disk format as any
struct.

Thirty-two bits of hash means a directory reaches even odds of \emph{some}
collision at roughly 65{,}000 entries. That is not exotic, and it is why the
name is in the payload at all: the key cannot distinguish colliding names, so
\fnname{lookup} compares the stored name after finding the key.

\decide{An item currently holds one entry, so a colliding pair cannot both be
stored. \fnname{bitter\_dir\_item} does not foreclose the fix --- \field{name\_len}
gives the stride and \field{bitter\_item.size} gives the end, so an item can hold
a sequence --- but nothing implements it. Harmless while nothing creates files.
See \code{docs/LOG.md}.}

\section{The inode item}

\bytes{112}, and the hardest struct in the project to change: every inode on
disk carries it.

It contains no inode number. The key's objectid \emph{is} the inode number, and
storing it twice would create two values that can disagree --- a corruption
\fnname{fsck} would then have to check for. The same rule explains why
\code{bitter\_root\_item} has no objectid: the payload never repeats what the
key already says.

Two fields are conditionally valid and look universal. \field{rdev} means
something only for \code{S\_ISCHR} and \code{S\_ISBLK}; \field{next\_dir\_index}
only for directories. Both are zero otherwise, and both are read without
checking \field{mode} by anyone who has not been told.

\paragraph{And there are no block pointers.} A classic Unix inode \emph{is} its
block map: direct pointers, then single, double and triple indirect. There is
nothing like that here. A file's contents are separate
\keytriple{ino}{EXTENT\_DATA}{file\_offset} items, and the B-tree supplies the
indirection at whatever depth the file's size demands --- so there is no
special case to write and no limit to grow out of. \field{size} and
\field{nbytes} are \emph{summaries} of that map rather than the map, and the
difference between them is exactly the holes.

\section{Where a file's bytes are}

Four numbers, in three coordinate systems, and none derivable from the others:

\begin{quote}
\field{disk\_bytenr} + \field{disk\_num\_bytes} --- the \emph{physical} extent\\
\field{offset} + \field{num\_bytes} --- this file's \emph{window} into it\\
the key's offset --- where that window lands \emph{in the file}
\end{quote}

\field{disk\_num\_bytes} must describe the whole extent rather than this file's
share, because with \field{disk\_bytenr} it is the extent tree's key
\keytriple{bytenr}{EXTENT\_ITEM}{length}; store the window and refcounting
searches for a key that does not exist.

\field{offset} is the one that looks redundant. Overwrite \kib{4} in the middle
of a \mib{1} file: copy-on-write puts the new bytes elsewhere, and the file then
needs three items, the first and third both naming the original extent at
different offsets into it. Without the field the only way to express that is to
split the extent physically --- a megabyte of I/O to change four kilobytes,
which defeats the point of CoW entirely. It is the same mechanism that makes
reflink possible.

\subsection{Reading one}

\fnname{bitterfs\_get\_block} answers one question: where does file block $n$
live? The answer has two forms and neither is the return value --- a mapped
buffer head, or an untouched one meaning \emph{hole}. The generic code
zero-fills for the second, so the sparse case the skeleton called this phase's
interesting one requires no code at all. It is the absence of an answer.

The lookup is by \emph{containment}, not identity: the item covering an offset
is the one with the largest key offset at or below it. A miss therefore steps
back one slot --- and needs no backward walk, because
\fnname{btree\_search}'s descent already steps to the child \emph{before} the
first separator greater than the key, landing in the leaf that holds the
predecessor even when that is a different leaf from the successor. Forward scans
need \fnname{btree\_next\_leaf}; containment lookups do not.

Everything above it is borrowed. \fnname{mpage\_read\_folio} owns the folio
locking, the bio submission and the zero-filling; \fnname{bitterfs\_read\_folio}
is one line.

\decide{A data block is the only thing on the device nothing verifies --- no
header, no checksum, no fsid. \fnname{get\_block} bounds and aligns the address
it computes, which catches a wild value and nothing subtler. Data checksums are
what \itemtype{EXTENT\_CSUM} is reserved for. See \code{docs/LOG.md}.}

\section{What the tests assert}

\code{qemu/guest.sh} grew from 15 checks to 36, and the ones with teeth are the
values read back out of \fnname{stat}: inode 256 and 257, link counts, modes,
and timestamps with nanoseconds intact. Each travelled from
\fnname{format\_initial\_fs\_tree} through the on-disk item, through
\fnname{read\_block}'s three checks, through byte-order conversion, into a
\code{struct inode}. A wrong \fnname{offsetof} or a missed accessor shows up as
a wrong number rather than as something that merely works.

Two are worth naming. \code{wc -c} returning exactly 20 proves \field{i\_size}
bounds the read and the \bytes{4076} of zero padding in the allocated block are
not returned. And a missing name being \emph{absent} rather than an error proves
\fnname{lookup} returns a negative dentry --- the distinction that decides
whether phase~7's \code{open(O\_CREAT)} can work at all.
